\section*{Chapter \ref{ch:b2:reaction-free-energy} --- Entropy and Free Energy of Reaction}

\begin{solution}{exo:b2:reaction-free-energy:1}
(a) Three moles of gas become liquid: $\Delta_r S^\circ < 0$. (b) A gas is
made: $> 0$. (c) One gas molecule gives two: $> 0$. (d) Ions in solution
form a crystal: $< 0$. (e) Melting: $> 0$. (f) $\Delta\nu_{\text{gas}} = 0$:
small, and here positive (\qty{+20}{J/(K.mol)} from the tables, \ce{HCl}
being less symmetrical than \ce{H2} and \ce{Cl2}).
% ledger: s0:HCl_g, s0:H2_g, s0:Cl2_g
\end{solution}

\begin{solution}{exo:b2:reaction-free-energy:2}
$\Delta_r S^\circ = 2(192.77) - 191.61 - 3(130.68) = \qty{-198.1}{J/(K.mol)}$;
$\Delta_r G^\circ = -91.9 - 298.15 \times (-0.1981) = \qty{-32.8}{kJ/mol}$:
\omterm{def:b2:reaction-free-energy:exergonic}{exergonic} at \qty{298}{K}, though the entropy term (\qty{+59.1}{kJ/mol})
works against it.
\end{solution}

\begin{solution}{exo:b2:reaction-free-energy:3}
$\Delta_r S^\circ = 69.95 - 130.68 - \tfrac12(205.15) = \qty{-163.3}{J/(K.mol)}$;
$\Delta_r G^\circ = -285.83 + 298.15 \times 0.1633 = \qty{-237.14}{kJ/mol}$.
It is the formation reaction of liquid water (one mole from the elements in
their reference states), so its \omterm{def:b2:reaction-free-energy:reaction-gibbs}{standard reaction Gibbs energy} is by
definition $\Delta_f G^\circ(\ce{H2O}, \text{l})$; the tables give the same
value.
\end{solution}

\begin{solution}{exo:b2:reaction-free-energy:4}
About \qty{42}{J/(K.mol)} at \qty{298}{K} and \qty{87}{J/(K.mol)} at
\qty{1000}{K} (liquid). Jumps: \qty{10.6}{J/(K.mol)} at \qty{692.7}{K} and
\qty{97.7}{J/(K.mol)} at \qty{1180.2}{K}. $\Delta_{\text{fus}}H = 10.6 \times
692.7 = \qty{7.3}{kJ/mol}$ and $\Delta_{\text{vap}}H = 97.7 \times 1180.2 =
\qty{115}{kJ/mol}$.
\end{solution}

\begin{solution}{exo:b2:reaction-free-energy:5}
$\Delta_r H^\circ = \qty{44.00}{kJ/mol}$, $\Delta_r S^\circ = 188.84 - 69.95 =
\qty{118.9}{J/(K.mol)}$, $\Delta_r G^\circ = 44.00 - 298.15 \times 0.1189 =
\qty{+8.55}{kJ/mol}$. At \qty{298}{K} the change is not reversible (liquid
water and vapour at \qty{1}{bar} are not in equilibrium), so $\Delta S \ne
Q/T$: $\Delta_r H^\circ/T = \qty{147.6}{J/(K.mol)}$. The positive
$\Delta_r G^\circ$ means that at \qty{298}{K} water does not evaporate into
an atmosphere of pure vapour at \qty{1}{bar}: the vapour condenses. The two
are in equilibrium where $\Delta_r G^\circ = 0$: $T = 44\,004/118.9 =
\qty{370}{K}$, within \qty{3}{K} of the normal boiling point
(\qty{373}{K}): the \omterm{def:b2:reaction-free-energy:ellingham-approximation}{Ellingham approximation} is good over \qty{75}{K}.
\end{solution}

\begin{solution}{exo:b2:reaction-free-energy:6}
$T_i = 57\,100/175.7 = \qty{325}{K}$. Below it $\Delta_r G^\circ > 0$ and
colourless \ce{N2O4} is favoured; cooling the tube shifts the gas towards
\ce{N2O4} and the brown colour of \ce{NO2} fades (the quantitative link with
the composition is in \cref{ch:b2:equilibrium-shifts}).
\end{solution}

\begin{solution}{exo:b2:reaction-free-energy:7}
$\Delta_r S^\circ = -\dd\Delta_r G^\circ/\dd T = -(-4.0 + 12.0)/100 =
\qty{-0.080}{kJ/(K.mol)} = \qty{-80}{J/(K.mol)}$ and $\Delta_r H^\circ =
\Delta_r G^\circ + T\Delta_r S^\circ = -12.0 + 300 \times (-0.080) =
\qty{-36.0}{kJ/mol}$. Check: $\Delta_r G^\circ/T$ goes from $-0.0400$ to
$\qty{-0.0100}{kJ/(K.mol)}$, slope $3.00 \times 10^{-4}$ per kelvin, equal to
$-\Delta_r H^\circ/T^2 = 36.0/346^2 = 3.00 \times 10^{-4}$ at the geometric
mean temperature $\sqrt{300 \times 400} = \qty{346}{K}$.
\end{solution}

\begin{solution}{exo:b2:reaction-free-energy:8}
$\Delta_r S^\circ = 186.25 - 5.74 - 2(130.68) = \qty{-80.8}{J/(K.mol)}$,
$\Delta_r G^\circ = -74.87 + 298.15 \times 0.0808 = \qty{-50.8}{kJ/mol}$,
the tabulated $\Delta_f G^\circ$ of methane. Both terms having the same
sign, methane becomes unstable above $T_i = 74\,870/80.8 \approx
\qty{926}{K}$ (\omterm{def:b2:reaction-free-energy:ellingham-approximation}{Ellingham approximation}): hot methane cracks to carbon and
hydrogen, a route to both.
\end{solution}

\begin{solution}{exo:b2:reaction-free-energy:9}
Substitution: \ce{CH3CHBrCH3 + OH- -> CH3CH(OH)CH3 + Br-}; elimination:
\ce{CH3CHBrCH3 + OH- -> CH2=CHCH3 + H2O + Br-}. The difference of the two
standard reaction Gibbs energies is $\Delta(\Delta_r G^\circ) = 40 - T \times
0.090$, negative above $40\,000/90 \approx \qty{440}{K}$. Elimination makes
one more molecule than substitution: its entropy is larger, and the entropy
term grows in proportion to $T$.
\end{solution}

\begin{solution}{exo:b2:reaction-free-energy:10}
Each of the $N_A$ molecules has 2 orientations: $W = 2^{N_A}$, so $S =
k_B\ln 2^{N_A} = R\ln 2 = \qty{5.76}{J/(K.mol)}$. The crystal is not
perfectly ordered: at low temperature the molecules are frozen in their
random orientations and cannot reach the single ordered arrangement the third
law assumes.
\end{solution}

\begin{solution}{exo:b2:reaction-free-energy:11}
Integrate $\dd\Delta_r S^\circ/\dd T = c/T$: $\Delta_r S^\circ(T) = \Delta_r
S^\circ(T_0) + c\ln(T/T_0)$; Kirchhoff gives $\Delta_r H^\circ(T) = \Delta_r
H^\circ(T_0) + c(T - T_0)$; subtract $T\Delta_r S^\circ(T)$. For limestone at
\qty{1100}{K}: Ellingham $178.32 - 1100 \times 0.16064 = \qty{1.62}{kJ/mol}$;
correction $c[T - T_0 - T\ln(T/T_0)] = -0.00195 \times (801.85 - 1436.3) =
\qty{+1.24}{kJ/mol}$: $\Delta_r G^\circ = \qty{2.86}{kJ/mol}$. The error,
\qty{1.2}{kJ/mol}, is less than one per cent of $\Delta_r H^\circ$, but near
the \omterm{def:b2:reaction-free-energy:inversion-temperature}{inversion temperature} it decides the sign.
\end{solution}

\begin{solution}{exo:b2:reaction-free-energy:12}
(a) $\Delta_r H^\circ = -763.2 + 944.0 = \qty{+180.8}{kJ/mol}$, $\Delta_r
S^\circ = 353.2 + 205.15 - 50.62 - 2(223.08) = \qty{+61.6}{J/(K.mol)}$;
$\Delta_r G^\circ(1000) = 180.8 - 61.6 = \qty{+119.2}{kJ/mol}$: \omterm{def:b2:reaction-free-energy:exergonic}{endergonic},
no chlorination. (b) For \ce{C + O2 -> CO2}, $\Delta_r G^\circ(1000) =
-393.5 - 2.9 = \qty{-396.4}{kJ/mol}$. The sum, \ce{TiO2(s) + 2Cl2(g) + C(s)
-> TiCl4(g) + CO2(g)}, has $\Delta_r G^\circ = \qty{-277.2}{kJ/mol}$: carbon
consumes the dioxygen made by the first reaction, and the coupled reaction is
strongly \omterm{def:b2:reaction-free-energy:exergonic}{exergonic}.
\end{solution}

\begin{solution}{pb:b2:reaction-free-energy:1}
\textbf{1.} \ce{CaCO3(s) -> CaO(s) + CO2(g)}.
\textbf{2.} $\Delta_r H^\circ = -635.09 - 393.51 + 1206.92 =
\qty{+178.3}{kJ/mol}$: \omterm{def:b2:reaction-enthalpy:exothermic}{endothermic}.
\textbf{3.} $\Delta_r S^\circ = 39.75 + 213.79 - 92.9 = \qty{+160.6}{J/(K.mol)}$,
positive because a gas is formed from a solid.
\textbf{4.} $\Delta_r G^\circ = 178.32 - 298.15 \times 0.16064 =
\qty{+130.4}{kJ/mol}$.
\textbf{5.} Yes: $\Delta_r G < 0$ is needed for decomposition, and here it
is strongly positive; the carbon dioxide of the air, far below \qty{1}{bar},
lowers $\Delta_r G$ but by far less than \qty{130}{kJ/mol} (next chapter).
\textbf{6.} $\Delta_r C^\circ_p = 42.80 + 37.13 - 81.88 =
\qty{-1.95}{J/(K.mol)}$, tiny: $\Delta_r H^\circ$ and $\Delta_r S^\circ$
hardly change with $T$.
\textbf{7.} $\Delta_r G^\circ(T) = 178.32 - 0.16064\,T$ (\unit{kJ/mol}).
\textbf{8.} \qty{+17.7}{kJ/mol} at \qty{1000}{K}; \qty{-30.5}{kJ/mol} at
\qty{1300}{K}.
\textbf{9.} $T_i = 178\,320/160.64 = \qty{1110}{K}$.
\textbf{10.} $\Delta_r G^\circ/T = 178.32/T - 0.16064$, whose derivative is
$-178.32/T^2 = -\Delta_r H^\circ/T^2$.
\textbf{11.} At $T_i$, $c[T_i - T_0 - T_i\ln(T_i/T_0)] = -0.00195 \times
(811.9 - 1459.4) = \qty{+1.26}{kJ/mol}$; $\Delta_r G^\circ$ reaches zero
$1260/160.6 \approx \qty{8}{K}$ higher, at about \qty{1118}{K}: less than
one per cent.
\textbf{12.} At \qty{1000}{K}, $\Delta_r G > 0$: the decomposition cannot
proceed (lime would absorb carbon dioxide). At \qty{1300}{K}, $\Delta_r G <
0$: limestone decomposes.
\textbf{13.} $\Delta_r G$ becomes positive: lime takes up the carbon dioxide
and turns back into carbonate.
\textbf{14.} $\delta S_{\text{c}} = -\Delta_r G\,\dd\xi/T = 30\,506/1300 =
\qty{23.5}{J/K}$ per mole.
\textbf{15.} Yes: \cref{ch:b2:chemical-potential} will show that
$\Delta_r G = \Delta_r G^\circ + RT\ln(p_{\ce{CO2}}/p^\circ)$, lowered
when the pressure of carbon dioxide is below \qty{1}{bar}, so that the
decomposition starts below $T_i$.
\textbf{16.} $n = 10^6/56.08 = \qty{1.783e4}{mol}$.
\textbf{17.} $1.783 \times 10^4 \times 178.32 = \qty{3.18e6}{kJ}$, that is
\qty{3.18}{GJ} per tonne.
\textbf{18.} $1.783 \times 10^4 \times 44.01 = \qty{785}{kg}$ of \ce{CO2}.
\textbf{19.} Heat to supply $3.18 \times 10^6/0.50 = \qty{6.36e6}{kJ}$, so
$6.36 \times 10^6/802.3 = \qty{7.93e3}{mol}$ of methane, \qty{127}{kg}.
\textbf{20.} $7.93 \times 10^3 \times 44.01 = \qty{349}{kg}$ of \ce{CO2}.
\textbf{21.} $785 + 349 = \qty{1134}{kg}$ per tonne of lime, 69\,\% of it
from the limestone.
\textbf{22.} The limestone's carbon dioxide is set by the stoichiometry,
whatever heats the kiln; only the fuel's share can be cut.
\textbf{23.} Limestone decomposes under \qty{1}{bar} of carbon dioxide above
its \omterm{def:b2:reaction-free-energy:inversion-temperature}{inversion temperature}, $\boldsymbol{T_i \approx \qty{1.11e3}{K}}$
(about \qty{840}{\celsius}).
\end{solution}
